
% !TeX root = ../main.tex

\chapter*{Lista de acrónimos y siglas}
\addcontentsline{toc}{chapter}{Lista de acrónimos y siglas}

\begin{description}
	\item[BH] Benjamini--Hochberg. Procedimiento de control de la tasa de falsos
	descubrimientos en contrastes múltiples.
	
	\item[CEO] \textit{Chief Executive Officer}. Máximo responsable ejecutivo de
	una empresa.
	
	\item[EPS] \textit{Earnings Per Share}. Beneficio por acción.
	
	\item[FDR] \textit{False Discovery Rate}. Proporción esperada de falsos
	positivos entre los contrastes declarados significativos.
	
	\item[ICC] Coeficiente de correlación intraclase. En este trabajo se emplea
	la forma ICC(1,1), de efectos aleatorios de una vía y medida individual.
	
	\item[LLM] \textit{Large Language Model}. Modelo de lenguaje de gran tamaño.
	
	\item[MAS] \textit{Multi-Agent System}. Sistema multiagente.
	
	\item[MDE] Efecto mínimo detectable. Magnitud mínima que un diseño puede
	detectar con la potencia estadística disponible.
	
	\item[P/E] \textit{Price-to-Earnings ratio}. Relación entre el precio de la
	acción y el beneficio por acción.
	
	\item[vLLM] Motor de inferencia para modelos de lenguaje, empleado para
	servir los modelos abiertos de la campaña experimental.
	
\end{description}

\chapter*{Glosario de términos}
\addcontentsline{toc}{chapter}{Glosario de términos}

\begin{description}
	\item[Agente ciego] Agente del comité cuyo \textit{prompt} omite el atributo
	sensible, de modo que su entrada es idéntica en las variantes de una misma
	pareja. Actúa como \textit{null} interno del diseño.
	
	\item[Alias] Identificador neutro (\textit{Company 1} a \textit{Company 4})
	con el que se presentan las empresas a los agentes, en sustitución de su
	nombre real y su \textit{ticker}, que se conservan únicamente en el registro
	experimental.
	
	\item[Arquetipo] Combinación de un sector económico y un perfil financiero
	que define una empresa sujeto y su cesta asociada. El diseño emplea 21
	arquetipos.
	
	\item[Atributo sensible] Característica de la empresa sujeto que se manipula
	en el contrafactual sin alterar la información relevante para la decisión.
	Este trabajo estudia dos: el género de la persona que ocupa el puesto de CEO
	y la localización de la sede.
	
	\item[Cesta] Conjunto de cuatro empresas del mismo sector que se presenta al
	comité en una ejecución: una empresa sujeto y tres empresas de contexto.
	
	\item[Composición] Combinación de modelos de lenguaje asignada a los tres
	agentes del comité. Es homogénea cuando los tres comparten el mismo modelo y
	heterogénea cuando pertenecen a familias distintas.
	
	\item[Contrafactual pareado] Comparación entre dos evaluaciones de una misma
	empresa que solo difieren en el valor del atributo sensible, manteniendo
	constante toda la información relevante para la decisión. Constituye el
	principio de comparación sobre el que se construye el diseño experimental.
	
	\item[Empresa de contexto] Cada una de las tres empresas que acompañan a la
	empresa sujeto dentro de una cesta y que se mantienen constantes entre las
	variantes de una pareja.
	
	\item[Empresa sujeto] Empresa de la cesta sobre la que se aplica la
	manipulación contrafactual y sobre la que se mide el \textit{gap} de
	asignación.
	
	\item[\textit{Gap}] Diferencia entre la cantidad asignada a la empresa sujeto
	en una variante y la asignada en su condición de control, expresada en euros.
	Un valor negativo indica una penalización de la variante.
	
	\item[Génesis] Turno inicial de la ejecución ($t=0$), en el que cada agente
	decide de forma independiente antes de observar las respuestas del resto del
	comité.
	
	\item[Nivel de instrucción] Grado de especificación del rol asignado a los
	agentes. El diseño define tres niveles: genérico
	(\textit{level\_0\_neutral}), profesional (\textit{level\_1\_professional})
	e identitario (\textit{level\_2\_identity}).
	
	\item[\textit{Null} externo] Referencia empírica del ruido del sistema
	obtenida de la condición placebo, que emplea un conjunto de cestas distinto
	del brazo contrafactual.
	
	\item[\textit{Null} interno] Referencia empírica del ruido del sistema
	obtenida del agente ciego, que opera sobre las mismas cestas y las mismas
	parejas que el brazo contrafactual.
	
	\item[Pareja] Conjunto de variantes generadas a partir de un mismo arquetipo,
	que comparten identificador y orden de presentación de las empresas.
	
	\item[Perfil financiero] Situación financiera empleada para estratificar la
	muestra de empresas sujeto: \textit{growth}, \textit{value} o
	\textit{distressed}.
	
	\item[Placebo] Condición de control formada por parejas de cestas gemelas
	idénticas (\textit{placebo\_a} y \textit{placebo\_b}), sin manipulación de
	ningún atributo sensible.
	
	\item[Ratio de varianzas] Cociente entre la varianza del \textit{gap} de los
	agentes visibles y la del agente ciego en génesis. Constituye el desenlace
	principal del análisis.
	
	\item[Repetición] Cada una de las tres ejecuciones de una misma condición
	experimental, identificadas mediante las semillas 42, 123 y 456.
	
	\item[\textit{Snapshot}] Descarga única y congelada de la información
	financiera empleada, obtenida el 12 de mayo de 2026, sobre la que se ejecutan
	de forma determinista todos los experimentos.
	
	\item[Vacuna] Instrucción de mitigación añadida al \textit{prompt} de sistema
	con el objetivo de reducir el efecto del atributo sensible sobre la decisión.
	
	\item[Variante] Cada una de las condiciones generadas sobre una misma empresa
	sujeto: control, variante de género y variante geográfica.
\end{description}